\section{Construcción proyectiva de la recta real y estructura conjunta del continuo}
\label{sec:construccion-proyectiva-recta}

\subsection{Configuraciones finitas y sistema de restricciones}

Para cada profundidad \(k\geq0\), sea \(\mathsf A_k\) el conjunto finito de
prefijos celulares de APP. Una emisión admisible
\(\varepsilon:a\to a'\) determina, para
\(\diamond\in\{\Sigma,\Pi\}\), la división exacta
\begin{equation}
 \delta^\diamond(\varepsilon)
 =\operatorname{res}^\diamond(\varepsilon)
  +9\operatorname{quo}^\diamond(\varepsilon),
 \qquad
 \operatorname{res}^\diamond(\varepsilon)\in\{0,\ldots,8\},
 \qquad
 \operatorname{quo}^\diamond(\varepsilon)\in\mathbb Z.
 \label{eq:division-app-configuracion-nuclear}
\end{equation}
El levantamiento TRIT proporciona
\begin{equation}
 \Delta(\varepsilon)
 =\operatorname{or}(\varepsilon)+3\kappa(\varepsilon),
 \qquad
 \operatorname{or}(\varepsilon)\in\{-1,0,+1\},
 \qquad
 \kappa(\varepsilon)\in\mathbb Z.
 \label{eq:division-trit-configuracion-nuclear}
\end{equation}
La coordenada de fase y memoria pertenece a
\(\mathbb Z_{\geq0}\times\mathbb Z/9\mathbb Z\). Con
\(\bar r\in\{0,\ldots,8\}\), el odómetro nonádico es
\begin{equation}
 \sigma_9(w,r)
 =\left(
 w+\left\lfloor\frac{\bar r+1}{9}\right\rfloor,
 r+[1]_9\right),
 \qquad
 \sigma_9^9(w,r)=(w+1,r).
 \label{eq:odometro-configuraciones}
\end{equation}

\begin{definition}[Configuración genealógica completa]
Sea \(\gamma=\varepsilon_1\cdots\varepsilon_k\) una ruta admisible. Su
configuración de nivel \(k\) es
\begin{equation}
\begin{aligned}
 \widehat x_k(\gamma)=\bigl(
 &a_{\leq k};\operatorname{Sh}_{\leq k};
 \mathbf{res}_k,\mathbf{quo}_k;
 \mathbf{or}_k,\boldsymbol\kappa_k;
 \mathsf Q_k,g_k,q_k^{(9)},\beta_k,m_k;\\
 &\partial\gamma,\gamma,\operatorname{Reg}(\gamma)
 \bigr),
 \label{eq:configuracion-genealogica-nuclear}
\end{aligned}
\end{equation}
donde \(\operatorname{Sh}_j=(\Sigma_j,\Pi_j)\),
\(\mathsf Q_k=(q_k^{(9)},g_k)\),
\(\beta_k=[q_k^{(9)}]_{12}\), y
\begin{equation}
 \operatorname{Reg}(\gamma)
 =\bigl(\lambda(\varepsilon_1),\ldots,
 \lambda(\varepsilon_k)\bigr)
 \label{eq:registro-ordenado-nuclear}
\end{equation}
conserva fuente, destino, residuos, cocientes, orientación, acarreo,
frontera y fases anterior y posterior de cada transición.
\end{definition}

Sea
\begin{equation}
 \widehat{\mathcal X}_k
 =\{\widehat x_k(\gamma):|\gamma|=k\}.
 \label{eq:espacios-finitos-configuraciones}
\end{equation}
El truncamiento de la última emisión y la prolongación neutra definen
\begin{equation}
 \rho_{k+1,k}:\widehat{\mathcal X}_{k+1}
 \longrightarrow\widehat{\mathcal X}_k,
 \qquad
 \eta_k:\widehat{\mathcal X}_k
 \longrightarrow\widehat{\mathcal X}_{k+1},
 \qquad
 \rho_{k+1,k}\eta_k=\id.
 \label{eq:restriccion-seccion-configuraciones}
\end{equation}

\begin{proposition}[Finitud, no vacuidad y sobreyectividad]
\label{prop:finitud-sobreyectividad-configuraciones}
Si cada configuración posee un conjunto finito no vacío de emisiones
admisibles y la prolongación neutra está definida en todo nivel, entonces
\(\widehat{\mathcal X}_k\) es finito y no vacío para todo \(k\), y
\(\rho_{k+1,k}\) es sobreyectiva.
\end{proposition}

\begin{proof}
La afirmación se prueba por inducción. El nivel inicial es finito y no vacío.
La unión finita de los conjuntos finitos no vacíos de emisiones produce el
nivel siguiente. Para cada \(x\in\widehat{\mathcal X}_k\), la identidad
\(\rho_{k+1,k}(\eta_k(x))=x\) proporciona una preimagen explícita.
\end{proof}

\subsection{Límite inverso, cilindros y medida canónica}

El espacio de historias completas es
\begin{equation}
 \widehat\Omega
 =\varprojlim_k
 \bigl(\widehat{\mathcal X}_k,\rho_{k+1,k}\bigr)
 =\left\{(x_k)_{k\geq0}:
 \rho_{k+1,k}(x_{k+1})=x_k\right\}.
 \label{eq:limite-inverso-historias}
\end{equation}

\begin{theorem}[Compacidad y existencia de historias completas]
\label{thm:compacidad-historias-completas}
El espacio \(\widehat\Omega\) es compacto y no vacío. Cada elemento conserva,
a toda profundidad, las dos hojas APP, la orientación TRIT, el acarreo, la
fase, la memoria, la frontera, la ruta y el registro ordenado.
\end{theorem}

\begin{proof}
El producto \(\prod_{k\geq0}\widehat{\mathcal X}_k\), con cada factor
discreto, es compacto. Las ecuaciones
\(\rho_{k+1,k}(x_{k+1})=x_k\) definen un subconjunto cerrado. La
sobreyectividad de los enlaces garantiza la propiedad de intersección finita
de los cilindros compatibles; la compacidad proporciona un elemento de su
intersección. Las aplicaciones de restricción conservan por definición cada
campo del prefijo.
\end{proof}

Para \(x_k\in\widehat{\mathcal X}_k\), sea
\(\operatorname{Em}(x_k)\) el conjunto de sus emisiones admisibles y sea
\(d(x_k)=\#\operatorname{Em}(x_k)\). La probabilidad local uniforme
\begin{equation}
 p_{x_k}(\varepsilon)=\frac1{d(x_k)}
 \label{eq:probabilidad-local-configuracion}
\end{equation}
asigna a una historia finita
\(h=(\varepsilon_{k+1},\ldots,\varepsilon_N)\) el peso
\begin{equation}
 \mu_{k,N}(h)=\prod_{j=k}^{N-1}\frac1{d(x_j)}.
 \label{eq:peso-historia-finita}
\end{equation}
Los pesos son proyectivamente consistentes:
\begin{equation}
 \sum_h\mu_{k,N}(h)=1,
 \qquad
 (\pi_{N+1,N})_*\mu_{k,N+1}=\mu_{k,N}.
 \label{eq:consistencia-medida-proyectiva}
\end{equation}

\begin{theorem}[Medida genealógica]
\label{thm:medida-genealogica}
Existe una única medida de Borel \(\mu\) sobre \(\widehat\Omega\) cuyo valor
en cada cilindro de prefijo \([h]\) es \(\mu_{k,N}(h)\). Su soporte es
\(\widehat\Omega\). Toda simetría genealógica que preserve los grados de
emisión preserva \(\mu\).
\end{theorem}

\begin{proof}
La consistencia \eqref{eq:consistencia-medida-proyectiva} define una medida
finitamente aditiva sobre el álgebra de cilindros. La extensión de
Carathéodory determina una única medida de Borel. Todo cilindro no vacío
posee peso positivo, de donde se obtiene soporte pleno. Una simetría que
preserva los grados permuta historias de igual peso y, por tanto, preserva la
medida en el álgebra generadora y en su extensión boreliana.
\end{proof}

\subsection{Las cinco construcciones inseparables}

Cada emisión \(\varepsilon:x\to y\) determina simultáneamente las estructuras
\begin{align}
 \mathcal T_\varepsilon
 &=(U_\varepsilon,\sigma_x,\sigma_y,
 A_\varepsilon,\eta_\varepsilon),
 &&\text{transporte polar y telescopía espectral},
 \label{eq:componente-transporte-polar}\\
 \mathcal B_\varepsilon
 &=(b_\varepsilon,b_\varepsilon^+,b_\varepsilon^-;
 M^+,M^-,D,H,\kappa^{\rm can}),
 &&\text{balance firmado y memoria acumulativa},
 \label{eq:componente-balance-firmado}\\
 \mathcal I_\varepsilon
 &=(\operatorname{Cel}_{\TPK}(\varepsilon),q,B,W_c(q)),
 &&\text{incidencia celular y curvatura de acarreo},
 \label{eq:componente-incidencia-celular}\\
 \mathcal P_\varepsilon
 &=(I_y,J_y,\kappa_{y,N},\delta_y,C_{y,N},F_\varepsilon),
 &&\text{complejo de puertos y cono del retorno},
 \label{eq:componente-complejo-puertos}\\
 \mathcal H_\varepsilon
 &=\Bigl(\Psi_\varepsilon,K_\varepsilon,
 \operatorname{Det}(C_{y,N}),H_*(C_{y,N}),\notag\\
 &\hspace{2.2em}\operatorname{Smith}(C_{y,N}),
 \operatorname{Bock}(C_{y,N})\Bigr),
 &&\text{invariantes homológicos y determinantales}.
 \label{eq:componente-homologico-determinantal}
\end{align}
Las cinco componentes poseen la misma fuente, el mismo destino, la misma
ruta y el mismo registro. Sea \(\mathcal E_k\) el espacio finito de
configuraciones completas de profundidad \(k\), incluidas sus emisiones
supervivientes, y sea \(\mathcal Z_k\) el producto tipado de los cinco
codominios. El ensamblaje es
\begin{equation}
 \mathcal O_k:\mathcal E_k\longrightarrow\mathcal Z_k,
 \qquad
 \widehat x_k\longmapsto
 \bigl(\mathcal T_k,\mathcal B_k,\mathcal I_k,
       \mathcal P_k,\mathcal H_k\bigr).
 \label{eq:ensamblaje-cinco-construcciones}
\end{equation}
Los mapas de restricción satisfacen
\begin{equation}
 u^{\mathcal Z}_{\ell,k}\circ\mathcal O_\ell
 =\mathcal O_k\circ u^{\mathcal E}_{\ell,k},
 \qquad \ell\geq k.
 \label{eq:naturalidad-cinco-construcciones}
\end{equation}

\begin{theorem}[Construcción conjunta del continuo]
\label{thm:cinco-construcciones-inseparables}
Existe una única aplicación
\begin{equation}
 \mathcal O_\infty:
 \Cdisc:=\varprojlim_k\mathcal E_k
 \longrightarrow
 \mathcal Z_\infty:=\varprojlim_k\mathcal Z_k
 \label{eq:operador-conjunto-continuo}
\end{equation}
cuyas proyecciones finitas son las aplicaciones
\(\mathcal O_k\). Las realizaciones espectral, solenoidal, gauge,
cohomológica y determinantal son proyecciones tipadas de
\(\mathcal O_\infty\); comparten una única sección de dominio.
\end{theorem}

\begin{proof}
La igualdad \eqref{eq:naturalidad-cinco-construcciones} convierte
\((\mathcal O_k)_k\) en un cono compatible sobre el sistema
\((\mathcal Z_k)_k\). La propiedad universal del límite inverso determina
una única aplicación \(\mathcal O_\infty\). Cada realización se obtiene al
componerla con la proyección correspondiente del producto tipado. La
eliminación de una componente modifica el codominio y destruye la
reconstrucción conjunta de los invariantes que esa componente conserva.
\end{proof}

Esta inseparabilidad es estructural. Transporte, balance, incidencia,
cohomología y determinante no constituyen cinco mecanismos de generación.
Constituyen cinco coordenadas functoriales del mismo sistema de historias.
El constructor del continuo es la aplicación
\begin{equation}
 G_{\rm cont}:
 \varprojlim_N\operatorname{Surv}^{(\chi)}_N
 \longrightarrow\Cdisc,
 \label{eq:constructor-conjunto-continuo}
\end{equation}
y \(\mathcal O_\infty\) actúa sobre el objeto ya construido por
\(G_{\rm cont}\).

\subsection{Evaluación arquimediana y reconstrucción de
\texorpdfstring{\(\mathbb R\)}{R}}

Para \(L\geq0\) y \(N\in\mathbb Z\), sea
\begin{equation}
 J_{N,L}=
 \left[\frac{N}{3^L},\frac{N+1}{3^L}\right],
 \qquad
 \operatorname{diam}J_{N,L}=3^{-L}.
 \label{eq:cilindros-ternarios-cerrados}
\end{equation}
La etiqueta posicional adopta la convención semiabierta
\(\{x\in\mathbb R:N3^{-L}\leq x<(N+1)3^{-L}\}\), excepto en el extremo derecho de una carta
compacta. Esta convención asigna una sola expansión a cada frontera; el
intervalo cerrado \(J_{N,L}\) se emplea en el argumento de compacidad.

Un lector arquimediano \(L_{\rm arq}\) asigna a
\(\widehat x=(x_k)_k\in\widehat\Omega\) una familia encajada
\(J_k(\widehat x)\) tal que
\begin{equation}
 J_{k+1}(\widehat x)\subseteq J_k(\widehat x),
 \qquad
 \operatorname{diam}J_k(\widehat x)\longrightarrow0.
 \label{eq:encaje-contraccion-arquimediana}
\end{equation}
Se define
\begin{equation}
 \operatorname{Val}(\widehat x)
 =\text{el único elemento de }
 \bigcap_{k\geq0}J_k(\widehat x).
 \label{eq:evaluacion-arquimediana-definitiva}
\end{equation}

\begin{theorem}[Recta real como cociente de historias]
\label{thm:recta-real-cociente-historias}
Supóngase que, para cada \(m\geq1\), los cilindros compatibles cubren
coherentemente \(K_m=[-m,m]\). Entonces la evaluación induce un
homeomorfismo
\begin{equation}
 \widehat\Omega^{(m)}/\!\sim_{\mathbb R}
 \xrightarrow{\;\cong\;}K_m,
 \qquad
 \widehat x\sim_{\mathbb R}\widehat y
 \Longleftrightarrow
 \operatorname{Val}(\widehat x)
 =\operatorname{Val}(\widehat y).
 \label{eq:homeomorfismo-cociente-real}
\end{equation}
Las inclusiones \(K_m\hookrightarrow K_{m+1}\) producen
\begin{equation}
 \varinjlim_m
 \bigl(\widehat\Omega^{(m)}/\!\sim_{\mathbb R}\bigr)
 \cong\mathbb R.
 \label{eq:recta-real-agotamiento}
\end{equation}
\end{theorem}

\begin{proof}
La intersección de compactos encajados es no vacía, y la convergencia de los
diámetros prueba su unicidad. El cilindro de una profundidad suficientemente
grande controla uniformemente la variación de la evaluación, lo que prueba
continuidad. La cobertura coherente construye una historia compatible sobre
cada punto de \(K_m\), de modo que la evaluación es sobreyectiva. El
cociente identifica exactamente sus fibras. La aplicación inducida es una
biyección continua desde un compacto hacia el espacio de Hausdorff \(K_m\);
por tanto, es un homeomorfismo. El agotamiento por los compactos
\([-m,m]\) prueba \eqref{eq:recta-real-agotamiento}.
\end{proof}

\subsection{Cifra, valor, número, medida, escala y dimensión}

\begin{definition}[Cifra]
Fijado un lector posicional \(L\), la cifra de profundidad \(k\) es
\begin{equation}
 d_{L,k}(x_k)=\pi_{L,k}^{\rm vis}(x_k).
 \label{eq:definicion-cifra}
\end{equation}
Es una coordenada finita del estado \(x_k\).
\end{definition}

\begin{definition}[Número y presentación]
Un número HMT es una sección compatible
\begin{equation}
 \boldsymbol x=(x_k)_{k\geq0}\in
 \varprojlim_k\operatorname{Surv}_k,
 \qquad \rho_{k+1,k}(x_{k+1})=x_k.
 \label{eq:definicion-numero-hmt}
\end{equation}
Una presentación es el par \((\boldsymbol x,L)\) formado por la sección y un
lector cuyo dominio la contiene.
\end{definition}

\begin{definition}[Valor]
El valor de una presentación arquimediana es
\begin{equation}
 \operatorname{Val}_L(\boldsymbol x)
 =\bigcap_{k\geq0}J_{L,k}(\boldsymbol x)\in\mathbb R,
 \label{eq:definicion-valor}
\end{equation}
donde la igualdad identifica la intersección unipuntual con su elemento.
La fibra \(\operatorname{Val}_L^{-1}(r)\) contiene las genealogías que el
lector \(L\) evalúa mediante la coordenada \(r\).
\end{definition}

\begin{definition}[Medida y escala]
Una medida es la composición
\begin{equation}
 (x,a)\in\mathcal C
 \longmapsto\mathcal O(x,a)
 \longmapsto\chi_u(\mathcal O(x,a))\in\mathbb R,
 \label{eq:definicion-medida}
\end{equation}
donde \(x\) es un estado, \(a\) es un lector o aparato,
\(\mathcal C\) es la relación de compatibilidad,
\(\mathcal O\) es el observable tipado y \(\chi_u\) es una carta de unidad.
La escala es la profundidad \(k\) junto con el diámetro del cilindro
evaluado. Refinar la escala sustituye un cilindro por un subcilindro y
conserva el prefijo genealógico.
\end{definition}

La clasificación aritmética pertenece al valor. En una carta de base mil,
\begin{equation}
 r=a_0+\sum_{n\geq1}\frac{b_n}{1000^n},
 \qquad b_n\in\{0,\ldots,999\},
 \label{eq:expansion-base-mil-clasificacion}
\end{equation}
una expansión racional termina exactamente cuando el denominador reducido
divide alguna potencia de \(1000\); es finalmente periódica exactamente
para los racionales; y es aperiódica para los irracionales. La algebraicidad
se determina por la anulación de un polinomio no nulo de
\(\mathbb Q[X]\); la trascendencia es el complemento de la clausura
algebraica de \(\mathbb Q\). La clasificación genealógica pertenece a la
sección y distingue estabilización, periodicidad completa, monodromía con
avance de memoria y prolongación coinductiva aperiódica.

La dimensión se determina antes de su realización métrica. Sea \(G\) el
grafo orientado de estados visibles y sean
\(c_1,\ldots,c_d\in Z^1(G;\mathbb Z)\) cociclos cuyas clases generan un
submódulo libre de rango \(d\) en \(H^1(G;\mathbb Z)\). Para una ruta
\(\gamma\),
\begin{equation}
 Q_j(\gamma)=\sum_{e\in\gamma}c_j(e),
 \qquad
 \widetilde\gamma=
 \bigl(q_{\rm vis}(\gamma),Q_1(\gamma),\ldots,Q_d(\gamma)\bigr).
 \label{eq:levantamiento-cohomologico-dimension}
\end{equation}

\begin{definition}[Dimensión genealógica]
La dimensión genealógica de un sistema de transporte es el rango del
submódulo de \(H^1(G;\mathbb Z)\) generado por las clases de memoria
independientes que deben conservarse para reconstruir sus rutas.
\end{definition}

\begin{theorem}[Minimalidad de las coordenadas de memoria]
\label{thm:minimalidad-dimension-memoria}
Sea \(Q:Z_1(G;\mathbb Z)\to\mathbb Z^d\) la evaluación de los \(d\)
cociclos independientes. Si \(m\) contadores aditivos se obtienen mediante
\(\widehat Q=AQ\), con \(A:\mathbb Z^d\to\mathbb Z^m\), y existe
\(B:\mathbb Z^m\to\mathbb Z^d\) tal que \(Q=B\widehat Q\), entonces
\begin{equation}
 m\geq d.
 \label{eq:cota-rango-dimension}
\end{equation}
\end{theorem}

\begin{proof}
La independencia de las clases implica
\(\operatorname{rank}\operatorname{im}Q=d\). Las factorizaciones dan
\[
 d=\operatorname{rank}\operatorname{im}Q
 \leq\operatorname{rank}\operatorname{im}\widehat Q
 \leq m.
\]
\end{proof}

Una pantalla dimensional
\(\Sigma_n(x)=(K_n,M_n(x),\beta_n(x),\mathcal R_n(x))\) realiza esas
memorias mediante el cociente
\begin{equation}
 d_{\Sigma_n}(x)
 =\dim_{K_n}\frac{M_n(x)}{\mathcal R_n(x)}
 =g_n-\operatorname{rank}_{K_n}R_n(x),
 \label{eq:dimension-realizada-cociente}
\end{equation}
cuando \(M_n(x)\cong K_n^{g_n}\) y las columnas de \(R_n(x)\) generan las
relaciones. La dimensión realizada es el número de canales independientes
que sobreviven a las relaciones de frontera. La cadena
\begin{equation}
 \text{estado TPK}\longrightarrow\Sigma_n(x)
 \longrightarrow s_D\in\Gamma(L_D)
 \xrightarrow{\;\chi_u\;}[s_D]_u\in\mathbb R
 \label{eq:cadena-medicion-dimensional}
\end{equation}
separa estado, tipo dimensional, magnitud y coordenada numérica.

\subsection{Doble evaluación del estado genealógico}

Sea \(\mathfrak X\subseteq\widehat\Omega\) el espacio compacto de historias
provistas de un calibre excepcional \(\chi_{\rm exc}\). Sobre una misma
sección actúan
\begin{equation}
 \mathfrak R:\mathfrak X\longrightarrow\mathbb R,
 \qquad
 \mathfrak E_{\chi_{\rm exc}}:
 \mathfrak X\longrightarrow\mathbf{Exc}.
 \label{eq:dos-evaluaciones-mismo-dominio}
\end{equation}
La primera contrae cilindros hasta su valor; la segunda conserva incidencia,
polaridad, orientación, frontera y cámara reticular.

\begin{theorem}[Doble evaluación]
\label{thm:doble-evaluacion}
La aplicación
\begin{equation}
 \mathfrak D_{\chi_{\rm exc}}(x)
 =\bigl(\mathfrak R(x),\mathfrak E_{\chi_{\rm exc}}(x)\bigr)
 \label{eq:aplicacion-doble-evaluacion}
\end{equation}
es la evaluación conjunta de una única sección. Si la familia
\(\{\mathfrak R,\mathfrak E_{\chi_{\rm exc}}\}\) separa puntos de
\(\mathfrak X\), entonces \(\mathfrak D_{\chi_{\rm exc}}\) es inyectiva.
\end{theorem}

\begin{proof}
Ambas aplicaciones poseen el mismo dominio y conmutan con las restricciones
del sistema proyectivo. Si las dos componentes de
\(\mathfrak D_{\chi_{\rm exc}}(x)\) y
\(\mathfrak D_{\chi_{\rm exc}}(y)\) coinciden, la propiedad de separación
implica \(x=y\).
\end{proof}

La realización arquimediana y la excepcional son dos proyecciones
matemáticas del estado genealógico. La realización dimensional es un funtor
posterior hacia secciones de rectas dimensionales. La distinción de
codominios preserva la unidad causal y evita identificar valor, incidencia y
magnitud.

\subsection{Alcance cardinal y relación con ZFC}

Sea \(H\) un espacio polaco de historias con topología de cilindros. Una
realización analítica es un conjunto \(A\subseteq\mathbb R\) para el que
existen un boreliano
\(C\subseteq H\times\mathbb N^{\mathbb N}\) y un lector boreliano
\(F:H\to\mathbb R\) tales que
\begin{equation}
 A=\{F(x):\exists z\in\mathbb N^{\mathbb N},\ (x,z)\in C\}.
 \label{eq:publicacion-analitica-hmt}
\end{equation}

\begin{theorem}[Dicotomía cardinal de las realizaciones analíticas]
\label{thm:dicotomia-cardinal-analitica}
Toda realización \(A\) de la forma
\eqref{eq:publicacion-analitica-hmt} es numerable o posee cardinal
\(2^{\aleph_0}\).
\end{theorem}

\begin{proof}
El conjunto \(A\) es analítico como imagen boreliana de la proyección de un
boreliano en un espacio polaco. El teorema del conjunto perfecto establece
que todo analítico no numerable contiene una copia homeomorfa del espacio de
Cantor \(2^{\mathbb N}\). Por tanto,
\(2^{\aleph_0}\leq|A|\leq|\mathbb R|=2^{\aleph_0}\).
\end{proof}

El cardinal inicial relativo a esta categoría es
\begin{equation}
 \aleph_{1,\HMT}^{\rm an}
 :=\min\{|A|:A\text{ es una realización analítica no numerable}\},
 \qquad
 \aleph_{1,\HMT}^{\rm an}=2^{\aleph_0}.
 \label{eq:cardinal-analitico-hmt}
\end{equation}
Este símbolo designa el primer cardinal no numerable dentro de la clase de
realizaciones analíticas. El cardinal clásico \(\aleph_1\) es el primer
ordinal no numerable de un universo conjuntista. La igualdad entre ambos es
exactamente la hipótesis clásica del continuo.

Sea
\begin{equation}
 Q_{\HMT}:=\widehat\Omega/\!\sim_{\mathbb R}\cong\mathbb R.
 \label{eq:cociente-hmt-cardinal}
\end{equation}
En ZFC son equivalentes
\begin{equation}
\begin{aligned}
 2^{\aleph_0}=\aleph_1
 &\Longleftrightarrow Q_{\HMT}\hookrightarrow\omega_1\\
 &\Longleftrightarrow
 \exists\rho:\widehat\Omega\to\omega_1,\quad
 \rho(x)=\rho(y)\Longleftrightarrow
 \mathfrak R(x)=\mathfrak R(y)\\
 &\Longleftrightarrow
 Q_{\HMT}\text{ admite un buen orden de tipo }\omega_1.
\end{aligned}
 \label{eq:equivalencias-ch-hmt}
\end{equation}
La construcción topológica
\(Q_{\HMT}\cong\mathbb R\) y la existencia de un buen orden de tipo
\(\omega_1\) son propiedades distintas. La primera se demuestra mediante
cilindros y cocientes; la segunda depende del universo conjuntista.

La insuficiencia genealógica del formalismo extensional se expresa de forma
precisa: el elemento \(r\in\mathbb R\) determina una clase de equivalencia en
\(Q_{\HMT}\), mientras ZFC no asigna canónicamente a \(r\) una sección de
\(\widehat\Omega\), su orientación, su acarreo, su fase o su memoria. Esta
afirmación caracteriza una diferencia de estructura, no una contradicción
formal de los axiomas de ZFC.

La presentación finita de APP, TRIT y TPK puede codificarse por un real
\(a_{\HMT}\subseteq\omega\). Sea
\begin{equation}
 U_{\HMT}=L[a_{\HMT}],
 \qquad
 \mathbb R^{U_{\HMT}}=\mathcal P(\omega)\cap U_{\HMT}.
 \label{eq:universo-constructible-hmt}
\end{equation}

\begin{theorem}[Hipótesis del continuo en el envolvente constructible]
\label{thm:ch-envolvente-constructible}
Existe en \(U_{\HMT}\), definible a partir de \(a_{\HMT}\), una biyección
\begin{equation}
 b_{\HMT}:\omega_1^{U_{\HMT}}
 \xrightarrow{\;\sim\;}\mathbb R^{U_{\HMT}},
 \qquad
 U_{\HMT}\models 2^{\aleph_0}=\aleph_1.
 \label{eq:biyeccion-ch-envolvente}
\end{equation}
\end{theorem}

\begin{proof}
Escribamos \(a=a_{\HMT}\) y \(M=L[a]\). Por condensación relativizada, todo
real \(x\in\mathbb R^M\) aparece en algún nivel
\(L_{\alpha+1}[a]\) con \(\alpha<\omega_1^M\). Defínase
\[
 \rho_a(x)=\min\{\alpha<\omega_1^M:x\in L_{\alpha+1}[a]\},
 \qquad
 R_\alpha=\{x:\rho_a(x)=\alpha\}.
\]
Cada \(R_\alpha\) es numerable en \(M\). Sea
\(j_\alpha:R_\alpha\hookrightarrow\omega\) la inyección mínima para el buen
orden constructible relativo y sea
\[
 D=\{(\alpha,j_\alpha(x)):
 \alpha<\omega_1^M,\ x\in R_\alpha\}.
\]
La aplicación \(E:D\to\mathbb R^M\), dada por
\(E(\alpha,j_\alpha(x))=x\), es biyectiva. El orden lexicográfico de \(D\)
posee tipo a lo sumo
\(\omega\cdot\omega_1^M=\omega_1^M\); como \(D\) no es numerable, su tipo
es \(\omega_1^M\). La composición de \(E\) con el isomorfismo de orden
\(\omega_1^M\to D\) proporciona \(b_{\HMT}\).
\end{proof}

El alcance lógico queda determinado por
\begin{equation}
 U_{\HMT}\models\mathsf{CH}.
 \label{eq:alcance-interno-ch}
\end{equation}
Una extensión exterior puede contener reales adicionales y satisfacer
\(\neg\mathsf{CH}\). La construcción genealógica resuelve la producción
topológica y dimensional del continuo y decide CH en su envolvente
constructible; la decisión de CH en un universo exterior requiere especificar
ese universo. Esta separación conserva la exactitud de los tres resultados:
la recta real es un cociente de historias, las realizaciones analíticas
satisfacen la dicotomía del conjunto perfecto y el envolvente
\(L[a_{\HMT}]\) satisface CH.
